% Auto-generated by Step 12 -- do not edit manually
\begin{table}[H]
\centering
\caption{\textbf{Sensitivity of Bunching Estimates to Polynomial Degree:
         Selected Sector-Year Cells}}
\label{tab:app_robust_polynomial}
\begin{threeparttable}
\setlength{\tabcolsep}{12pt}
\small
\begin{tabular}{lcccc}
\toprule
 & Degree 5 & Degree 6 & Degree 7 & Degree 8 \\
Sector-Year & & & (Baseline, $p=7$) & \\
\midrule
\multicolumn{5}{l}{\textit{Panel A: Pre-Reform (2015, $z^* = 10M$ MNT)}} \\[2pt]
\quad Retail 2015          & 4.050 & 3.790 & 3.820 & 3.652 \\
\quad Food service 2015    & 2.411 & 2.180 & 2.161 & 1.994 \\
\quad Wholesale 2015       & 2.033 & 2.069 & 2.120 & 2.341 \\
\quad Manufacturing 2015   & 1.655 & 1.665 & 1.676 & 1.575 \\
[6pt]
\multicolumn{5}{l}{\textit{Panel B: Post-Reform (2022, $z^* = 50M$ MNT)}} \\[2pt]
\quad Retail 2022          & 1.881 & 1.890 & 1.919 & 1.839 \\
\quad Food service 2022    & 3.678 & 3.618 & 4.099 & 6.532 \\
\bottomrule
\end{tabular}
\begin{tablenotes}[para,flushleft]\footnotesize
\item \textit{Notes:} Exclusion window $[z^*_t - 2M,\, z^*_t + 30M)$ post-reform, $[z^*_t - 2M,\, z^*_t + 10M)$ pre-reform.
Degree 7 is the baseline specification.
Estimates are generally stable across degrees 5--8.
\textit{Source:} Author's calculations based on MTA CIT registry data.
\end{tablenotes}
\end{threeparttable}
\end{table}
